{
    \setlength{\parindent}{0em}
    \par {\zihao{4}\bfseries 一、题目：\Title}
    \par {\zihao{4}\bfseries 二、指导教师对文献综述、开题报告、外文翻译的具体要求：}
    \vspace{0.6em}
    {\zihao{5}
        \par \textbf{文献综述要求：}系统检索近三年CCF A类及arXiv前沿论文，围绕“自适应检索”、“多源知识融合”、“复杂推理”等主题，按技术范式分类对比至少8-10篇核心工作，总结各方法的优势与局限，最终提炼出当前研究存在的2-3个关键问题。

        \vspace{0.4em}
        \par \textbf{开题报告要求：}
        \par 1.\quad 界定“多源异构”、“复杂推理”等核心概念，并对痛点进行分析，说明本文研究的必要性。
        \par 2.\quad 明确研究内容，将核心问题拆解为 3-4 个具体的子任务，且子任务之间需有递进关系。
        \par 3.\quad 提供详细的技术路线，论证可行性，并归纳创新点。
        \par 4.\quad 编制切实可行的进度计划，说明交付成果内容。

        \vspace{0.4em}
        \par \textbf{外文翻译要求：}
        \par 1.\quad 选自 CCF A 类会议近两年发表的论文。
        \par 2.\quad 题目直接涉及 “RAG”与“Reasoning” 或 “Knowledge Graph”与“Retrieval” 的结合。
        \par 3.\quad 正文部分（不含参考文献、附录）不少于5000-6000英文单词，对应中文译文不少于 8,000 字。
        \par 4.\quad 要求专业术语、长难句及技术细节等内容翻译准确，关键术语在首次出现时标注英文原文。
    }


}

\mbox{} \vfill

\signature{指导教师（签名）}
% comment the line above and uncomment the line below if you want to set a signature with a specific date.
% \signaturewithdate{指导教师（签名）}{1897}{5}{21}
